\chapter{Osservabilità e metriche}
\label{chap:osservabilita}

\section{Pubblicazione delle metriche}

Le metriche operative possono tipicamente raggiungere chi le osserva in due modi: o è qualcun altro a 
interrogare periodicamente ogni servizio (\emph{pull}), oppure è ogni servizio a inviarle di propria 
iniziativa (\emph{push}). 
RAM-USB adotta il secondo: ogni servizio pubblica le proprie metriche una volta al minuto,
sul proprio topic MQTT dedicato verso il broker Mosquitto \cite{mosquitto-github}. 
La scelta deriva dal fatto che un modello \emph{pull} richiederebbe che Metrics-Collector apra una 
connessione verso ciascuno dei sei servizi, ampliando le regole di reachability della mesh discusse nel 
capitolo 7.
Pubblicando di propria iniziativa, sono invece i servizi ad aprire la connessione verso il broker.
Da lì, Metrics-Collector legge le metriche e le persiste in un database.

Le metriche operative non contengono dati degli utenti, ma i metadati del sistema vanno protetti con la 
stessa cura, per questo le comunicazioni avvengono all'interno della rete privata e protette dal mTLS, 
certificato dalla CA interna.

Ogni pubblicazione usa \emph{QoS 1}, (at least once) nella classificazione MQTT \cite{oasis-mqtt-311}: il
broker garantisce che il messaggio arrivi almeno una volta, al costo di un possibile duplicato. \emph{QoS 0},
(at most once), è escluso perché RAM-USB non può permettersi di perdere metriche in silenzio; \emph{QoS 2},
(exactly once), è escluso perché la garanzia in più che offre (l'assenza di duplicati) qui non serve
visto che un duplicato si traduce solo in due punti sovrapposti nello stesso istante, invisibili sul
grafico.

Ogni pubblicazione porta come corpo un payload di sei campi, riportato nel \autoref{lst:metrics-payload}.
I campi sono rispettivamente: l'identificativo del servizio che sta pubblicando, l'istante della 
pubblicazione, il numero di richieste gestite nell'ultimo minuto, il numero di errori rilevati nell'ultimo 
minuto, il tempo di risposta medio dell'ultimo minuto in millisecondi e il numero di connessioni che 
risultano aperte in quel momento. 

\Needspace{10\baselineskip}
\begin{lstlisting}[caption={Struttura del payload delle metriche, da \texttt{payload.go}
\cite{ramusb-code-metrics-payload}.}, label={lst:metrics-payload}]
type Payload struct {
	Service               string  `json:"service"`
	Timestamp             string  `json:"timestamp"`
	RequestCount          int64   `json:"request_count"`
	ErrorCount            int64   `json:"error_count"`
	AverageResponseTimeMs float64 `json:"average_response_time_ms"` 
	ActiveConnections     int64   `json:"active_connections"`
}
\end{lstlisting}

\section{Ricezione e controllo di coerenza}

Il messaggio pubblicato deve arrivare a un solo destinatario, e a nessun altro. Questa proprietà è garantita 
dall'ACL di Mosquitto che garantisce che ogni servizio possa solo scrivere sul proprio topic, mentre 
Metrics-Collector può leggere dai topic \texttt{metrics/\#} senza potervi scrivere.

Quell'ACL, però, non impedisce che un servizio compromesso, pur restando entro il proprio permesso di 
scrittura legittimo, pubblichi un payload il cui campo \texttt{service} dichiarato sia un nome diverso dal 
proprio.
Metrics-Collector rivalida quindi ogni messaggio in modo indipendente, derivando il servizio
atteso dal topic di arrivo e confrontandolo con quello dichiarato nel payload. 
Su qualunque discrepanza scarta il messaggio e lo registra nei log.

\section{Persistenza delle serie temporali}

Metrics-Collector persiste le metriche in TimescaleDB \cite{timescaledb-github}, un'estensione di 
PostgreSQL pensata per un carico fatto quasi solo di scritture ordinate nel tempo. 

La tabella \texttt{metrics} non ha una chiave primaria dato che una riga di metriche non ha un'identità 
naturale, inoltre, TimescaleDB \cite{timescaledb-github} non ne richiede una. È organizzata in chunk per
intervalli di tempo sulla colonna \texttt{time}: le due politiche in background del
\autoref{lst:metrics-policies} agiscono chunk per chunk, in base alla loro età. La prima elimina ogni chunk
più vecchio di trenta giorni; la seconda, dopo sette giorni, sposta i chunk in un formato compresso a
colonne, raggruppato per servizio e ordinato per tempo decrescente.

\Needspace{18\baselineskip}
\begin{lstlisting}[language=sql, caption={Politiche automatiche di ritenzione e compressione della tabella
\texttt{metrics}, da \texttt{000001\_create\_metrics\_table.up.sql} \cite{ramusb-code-metrics-schema}.},
label={lst:metrics-policies}]
SELECT create_hypertable('metrics', by_range('time'));

SELECT add_retention_policy('metrics', drop_after => INTERVAL '30 days');

ALTER TABLE metrics SET (
	timescaledb.enable_columnstore,
	timescaledb.segmentby = 'service',
	timescaledb.orderby = 'time DESC'
);
CALL add_columnstore_policy('metrics', after => INTERVAL '7 days');
\end{lstlisting}

\section{Visualizzazione delle metriche}

Un amministratore consulta lo stato di salute del sistema, senza che quella consultazione 
tocchi mai i dati degli utenti visto che per costruzione, le metriche non ne contengono. 
Grafana \cite{grafana-github} legge da TimescaleDB direttamente.

Per contattare TimescaleDB, oltre alla raggiungibilità sulla rete mesh e al certificato mTLS di Grafana, 
è richiesta anche una password, sullo stesso principio di difesa in profondità già visto per l'API 
amministrativa di Headscale nel capitolo 7.

La dashboard mostra tre pannelli, ciascuno una query diretta sulla tabella \texttt{metrics}: tempo medio di
risposta, throughput e connessioni attive. 

La \autoref{fig:metrics-flow} riassume l'intera pipeline appena descritta, dalla pubblicazione sul topic
MQTT alla dashboard.

\begin{figure}[H]
    \centering
    \includegraphics[width=\textwidth]{figures/11-operations-metrics-flow.pdf}
    \caption{Pipeline di osservabilità, dalla pubblicazione delle metriche alla dashboard di Grafana.}
    \label{fig:metrics-flow}
\end{figure}
